Autonomous agents operating in unknown environments must act before the
full structure of the world is known, revise their behavior as geometry
and risk are revealed, and do so without destabilizing previously learned
competencies. Existing approaches address parts of this challenge, but
often hold fixed one of the very quantities that must remain adaptable in
changing environments: the policy, the objective, the planner, or the
underlying decision dynamics.

This thesis studies that gap through the lens of \textbf{continual
Hamiltonian learning}, in which an agent's policy is identified with the
flow map of a structured phase-space energy whose landscape is
progressively reshaped by new geometric and risk information. Under this
view, adaptation is expressed not as relearning a flat action map, but as
updating the force-field coefficients, objectives, and correction laws
that govern the induced dynamics.

The thesis develops this viewpoint in two stages. \textbf{Stage~1}
establishes a geometry-only foundation for structured navigation under
local sensing, showing that phase-space energy shaping can support stable
goal-seeking, obstacle avoidance, and online correction without requiring
a globally complete map. \textbf{Stage~2} enlarges this learner to the
material-aware risk-sensitive setting by introducing semantic risk terms
into the stored energy, CVaR-based tail-risk optimization into the
objective, and a passivity- and coverage-guided hierarchy of updates
across segment, episode, and curriculum timescales. On a 20-episode high-risk subset of the DFC2018
Houston benchmark (oracle material maps, spatially held-out test region),
the resulting material-aware learner reduces cumulative
risk exposure by 56.4\% relative to the geometry-only baseline and by 17.4\%
over a material-blind Dijkstra reference, while maintaining comparable path
efficiency (Stage~2 path length is 1.1\% shorter than Stage~1).

A central thesis-level insight is that enriching the learner with a new
energy term also enlarges the associated force channels, coefficient
space, objective, gradient pathways, and update laws required to train
it. This makes the transition from geometry-only to material-aware
learning structurally necessary rather than an ad hoc design choice.
Taken together, the results support continual Hamiltonian learning as a
promising framework for progressively enriching decision dynamics while
preserving interpretability, stability, and adaptability in unknown
environments.